% vim: ts=2:sw=2:expandtab
\documentclass[11pt]{article}

\usepackage[T1]{fontenc}
\usepackage[utf8]{inputenc}
\usepackage{lmodern}
\usepackage{a4wide}
\usepackage{amsmath,amssymb,amsthm}
\usepackage{hyperref}

\theoremstyle{definition}
\newcommand{\Nat}{\mathbb{N}}
\newcommand{\NatZero}{\mathbb{N}_0}
\newcommand{\Real}{\mathbb{R}}
\newcommand{\Int}{\mathbb{Z}}
\newcommand{\Rat}{\mathbb{Q}}
\newcommand{\Coef}{\mathbb{K}}
\newcommand{\Dom}{\mathbb{D}}
\newcommand{\mset}[1]{\{\!\{#1\}\!\}}
\newcommand{\cvec}{\overline{c}}
\newcommand{\bfm}{\mathbf{m}}
\newcommand{\bfDelta}{\boldsymbol{\Delta}}
\newcommand{\tickto}[1]{\xrightarrow{\mathit{tick}(#1)}}
\newcommand{\DCPN}[1]{\texorpdfstring{\textup{DCPN}\ensuremath{_{#1}}}{DCPN#1}}

\newtheorem{definition}{Definition}
\newtheorem{remark}{Remark}

\title{Semantics of Coloured Petri Nets in SDCPN sense}
\author{Igor Konnov, Thomas Pani}
\date{\today}

\begin{document}

\maketitle

\begin{abstract}
Coloured Petri nets were introduced by Jensen~\cite{Jensen1981}; see Jensen and
Kristensen~\cite{JensenKristensen2009} for a modern presentation. DCPNs and SDCPNs were introduced by
Everdij and Blom~\cite{EverdijBlom2005,EverdijBlom2006}. Even though both carry the word
``colored'' in their name, the behavior of colored tokens differs in (S)DCPNs
from the original semantics of CPNs. The purpose of this document is to capture
the semantics of the CPN fragment of DCPNs. This is needed to reason about
the small model theorem and further translation to reactive modules.
\end{abstract}

\section{\DCPN{0}}

In this section, we define \DCPN{0}, the most minimalistic fragment of DCPNs
that still makes sense:

\begin{itemize}
  \item Colors are unified, that is, in every place, a token has its color
    from~$\Dom^k$ for a fixed~$k \ge 1$. The domain~$\Dom$ is either the
    reals~$\Real$ or the naturals~$\NatZero$.

  \item There are only ordinary arcs, no enabling or inhibitor arcs.

  \item There are no delay transitions.

  \item Every transition has at least one input arc.

  \item Firing is deterministic: a transition produces exactly one token along
    each of its output arcs, and the colours of the produced tokens are an
    affine function of the colours of the consumed tokens.

  \item Guarded transitions have their guards over linear integer arithmetic
    (LIA) if $\Dom = \NatZero$, and over linear real arithmetic (LRA) if
    $\Dom = \Real$.

  \item Token color functions are functions of the token colors. There is a
    global clock that advances in ticks, whose durations range over~$\Dom$.
    In every tick, the token colors change by one Euler step of the tick
    duration.
\end{itemize}

%%%%%%%%%%%
\subsection{Structure and syntax}

We specialize the definition of DCPNs by Everdij and
Blom~\cite[Sec.~3.1]{EverdijBlom2005} to this fragment. We keep their tuple
structure and notation: calligraphic letters denote the elements of the tuple,
whereas $P$, $T$, and $A$ denote individual places, transitions, and arcs,
respectively.

\begin{definition}[\DCPN{0}]\label{def:dcpn0}
  Fix $k \in \Nat$ to be the color dimension, different from~0, and fix the
  domain $\Dom \in \{\Real, \NatZero\}$.
  A \emph{\DCPN{0}} is a tuple
  $(\mathcal{P}, \mathcal{T}, \mathcal{A}, \mathcal{N}, \mathcal{S},
    \mathcal{C}, \mathcal{V}, \mathcal{G}, \mathcal{D}, \mathcal{F},
    \mathcal{I})$,
  together with the rules R0--R4 of~\cite{EverdijBlom2005}, whose elements
  are restricted as follows:

  \begin{enumerate}
    \item $\mathcal{P}$ is a finite set of places.

    \item $\mathcal{T}$ is a finite set of transitions with $\mathcal{T} \cap
      \mathcal{P} = \emptyset$. In general, it is partitioned into guard
      transitions $\mathcal{T}_G$, delay transitions $\mathcal{T}_D$, and
      immediate transitions $\mathcal{T}_I$. In this definition, $\mathcal{T}_D
      = \emptyset$. Hence, $\mathcal{T} = \mathcal{T}_G \cup \mathcal{T}_I$.

    \item $\mathcal{A}$ is a finite set of arcs with $\mathcal{A} \cap
      \mathcal{P} = \mathcal{A} \cap \mathcal{T} = \emptyset$.  In general, it
      is partitioned into ordinary arcs $\mathcal{A}_O$, enabling arcs
      $\mathcal{A}_E$, and inhibitor arcs $\mathcal{A}_I$. In this definition,
      $\mathcal{A}_E = \mathcal{A}_I = \emptyset$. Hence, $\mathcal{A} =
      \mathcal{A}_O$.

    \item $\mathcal{N} : \mathcal{A} \to
      (\mathcal{P} \times \mathcal{T}) \cup (\mathcal{T} \times \mathcal{P})$
      is a node function. For an arc $A$, we write $P(A)$ and $T(A)$ for the
      place and the transition of $\mathcal{N}(A)$, respectively. For a
      transition $T$, the sets $\mathcal{A}_{\mathit{in}}(T)$ and
      $\mathcal{A}_{\mathit{out}}(T)$ of its input and output arcs are defined
      as in~\cite{EverdijBlom2005}. Since all arcs are ordinary,
      $\mathcal{A}_{\mathit{in},O}(T) = \mathcal{A}_{\mathit{in},OE}(T)
        = \mathcal{A}_{\mathit{in}}(T)$. Note that
      $\mathcal{N}$ need not be injective, that is, a place and a transition
      may be connected by several arcs. In this definition, every transition
      has at least one input arc, that is,
      $\mathcal{A}_{\mathit{in}}(T) \ne \emptyset$ for every
      $T \in \mathcal{T}$.

    \item $\mathcal{S} = \{\Dom^k\}$ is the set of colour types. In this
      definition, there is only one colour type --- the set of
      vectors~$\Dom^k$.

    \item $\mathcal{C} : \mathcal{P} \to \mathcal{S}$ is the colour function.
      In our case, it is trivially defined as $\mathcal{C}(P) = \Dom^k$ for
      every $P \in \mathcal{P}$.  For a transition $T$ with input arcs
      $\mathcal{A}_{\mathit{in}}(T) = \{A_1, \dots, A_m\}$ and output arcs
      $\mathcal{A}_{\mathit{out}}(T) = \{A'_1, \dots, A'_n\}$, we write
      $\mathcal{C}(P(\mathcal{A}_{\mathit{in}}(T))) = \mathcal{C}(P(A_1))
      \times \dots \times \mathcal{C}(P(A_m)) = \Dom^{k \cdot m}$, and
      similarly $\mathcal{C}(P(\mathcal{A}_{\mathit{out}}(T)))
      = \Dom^{k \cdot n}$. While $m \ge 1$, a transition may have no output
      arcs, that is, $n = 0$, and then the latter product is empty.
      Unlike~\cite{EverdijBlom2005}, who set $\Real^0 = \emptyset$, we let
      $\Dom^0 = \{()\}$ be the singleton set containing the empty tuple.
      Hence, the firing measure $\mathcal{F}_T$ of a transition without
      output arcs is still a probability measure, and such a transition can
      fire.

    \item $\mathcal{V}$ is the set of token colour functions. For each
      place $P \in \mathcal{P}$, it contains a function
      $\mathcal{V}_P : \Dom^k \to \Dom_\pm^k$, which gives the rate of
      change of the colour of a token in~$P$. Here, $\Dom_\pm = \Real$ if
      $\Dom = \Real$, and $\Dom_\pm = \Int$ if $\Dom = \NatZero$, so that
      colours may also decrease. If $\Dom = \Real$, then $\mathcal{V}_P$ is
      locally Lipschitz continuous, as in~\cite{EverdijBlom2005}.

    \item $\mathcal{G}$ is the set of transition guards. For each
      $T \in \mathcal{T}_G$, it contains a guard
      $\mathcal{G}_T : \mathcal{C}(P(\mathcal{A}_{\mathit{in}}(T)))
        \to \{\mathit{True}, \mathit{False}\}$
      with $\mathcal{G}_T(\cvec) = \mathit{True}$ iff $\cvec \models \varphi_T$,
      where $\varphi_T$ is a quantifier-free formula over the colour
      components of the input tokens. If $\Dom = \NatZero$, then $\varphi_T$
      is a formula of linear integer arithmetic (LIA) with coefficients
      in~$\Int$, and if $\Dom = \Real$, then $\varphi_T$ is a formula of
      linear real arithmetic (LRA) with coefficients in~$\Rat$.\footnote{%
        Everdij and Blom~\cite{EverdijBlom2005} define
        $\mathcal{G}_T(\cvec) = \mathit{True}$ iff $\cvec \in \partial O_T$, that is,
        $\cvec$ lies on the boundary of an open set~$O_T$. Since the colours
        evolve continuously, a guard transition fires at the moment when its
        input colours reach this boundary. In contrast, our guard holds for
        every colour that satisfies~$\varphi_T$, which matches the predicate
        transitions of Petrinaut, which fire in every simulation step in which
        their predicate holds. Under a discretized clock, the boundary
        semantics would be fragile: a colour may jump over the boundary
        between two clock ticks, for example, from~$4$ to~$6$ over the
        boundary~$5$, and the transition would never fire.}

    \item $\mathcal{D} = \emptyset$, since there are no delay transitions.

    \item $\mathcal{F}$ is the set of firing measures. In general, for each
      $T \in \mathcal{T}$, it contains a map $\mathcal{F}_T$ from
      $\mathcal{C}(P(\mathcal{A}_{\mathit{in}}(T)))$ into the set of
      probability measures on
      $\{0, 1\}^{|\mathcal{A}_{\mathit{out}}(T)|}
        \times \mathcal{C}(P(\mathcal{A}_{\mathit{out}}(T)))$.
      A sample $(f, a)$ selects the output arcs that receive a token ($f$)
      and the colours of the produced tokens ($a$).
      In this definition, firing is deterministic. Let $T$ have $m$ input arcs
      and $n$ output arcs. Then there is an affine function
      $\pi_T(\cvec) = M_T\, \cvec + b_T$ with $M_T \in \Coef^{kn \times
      km}$ and $b_T \in \Coef^{kn}$, where $\Coef = \Int$ if
      $\Dom = \NatZero$, and $\Coef = \Rat$ if $\Dom = \Real$, such that
      $\pi_T(\cvec) \in \Dom^{kn}$ and
      $\mathcal{F}_T(\cvec) = \delta_{(\mathbf{1}, \pi_T(\cvec))}$ for all
      $\cvec \in \Dom^{km}$. Here, $\delta_x$ is the Dirac measure at~$x$, and
      $\mathbf{1} = (1, \dots, 1) \in \{0, 1\}^n$, that is, $T$ produces
      a token along each of its output arcs.

    \item $\mathcal{I} : \mathcal{P} \to \mathcal{C}(P)_{\mathit{ms}}$ is the
      initialisation function, where $\mathcal{C}(P)_{\mathit{ms}}$ is the set
      of all multisets over $\mathcal{C}(P) = \Dom^k$.
  \end{enumerate}
\end{definition}

The rules R0--R4 resolve competing enablings and are therefore part of the
execution semantics, which we define later. Since $\mathcal{T}_D = \emptyset$,
rule R0 reduces to immediate transitions having priority over guard
transitions. Likewise, the global clock only affects executions:
Definition~\ref{def:dcpn0} merely fixes that, while a token with colour~$c$
resides in a place~$P$, its colour changes at the rate~$\mathcal{V}_P(c)$.

\begin{remark}
  Everdij and Blom define $\mathcal{C}(P(\mathcal{A}_{\mathit{in},OE}(T)))$
  as a product over the \emph{places} connected to~$T$, whereas their
  execution consumes one token per \emph{arc}. We take the product over arcs,
  so that a transition with two arcs from the same place reads the colours of
  two distinct tokens.
\end{remark}

\begin{remark}[Transitions without input arcs]\label{rem:sources}
  DCPNs~\cite{EverdijBlom2005} and Petrinaut allow transitions without input
  arcs, often called \emph{sources}. Such transitions are meaningful only as
  delay transitions. A guard transition without input arcs never fires, since
  its guard is defined on $\Dom^0$ and has an empty boundary. An immediate
  transition without input arcs is always enabled, so it would fire infinitely
  often at the same time. In contrast, Everdij and Blom explicitly allow a
  delay transition without input arcs, whose enabling rate~$d_T$ is then a
  constant, that is, the transition fires as a Poisson process. Similarly, the
  Petrinaut examples use sources with stochastic rates to model exogenous
  arrivals, for example, customer demand or launched satellites. Since
  \DCPN{0} has no delay transitions, we require every transition to have at
  least one input arc. Tokens enter a \DCPN{0} only via the initial marking.
\end{remark}

%%%%%%%%%%%
\subsection{Small-step execution semantics}

Everdij and Blom define execution semantics over timestamps. In this sense,
they define executions somewhat similar to big-step semantics. In this section,
we define executions via individual transitions, closer to how it is done in
timed and hybrid automata. First, we define what it means to fire a single
transition, which we call a \emph{firing}. Firings operate on intermediate
configurations. A
maximal sequence of transitions is organized into a \emph{round}. A round
finishes with a clock tick.

Throughout this section, we fix a \DCPN{0} as in
Definition~\ref{def:dcpn0}. For every transition~$T$, we fix an enumeration
$A_1, \dots, A_m$ of its input arcs and an enumeration $A'_1, \dots, A'_n$ of
its output arcs, which we used to define the products
$\mathcal{C}(P(\mathcal{A}_{\mathit{in}}(T))) = \Dom^{km}$ and
$\mathcal{C}(P(\mathcal{A}_{\mathit{out}}(T))) = \Dom^{kn}$. For
$\cvec \in \Dom^{km}$, we write $\pi_T^j(\cvec) \in \Dom^k$ for the $j$-th block of
$k$ components of $\pi_T(\cvec)$, that is, the colour produced along~$A'_j$.

\paragraph{Markings.}
We write $\mset{\cdot}$ for multisets, $+$ and $-$ for their sum and
difference, and $\le$ for multiset inclusion. A \emph{token} is a pair
$(P, c) \in \mathcal{P} \times \Dom^k$ of a place and a colour, and a
\emph{marking} is a finite multiset of tokens. The \emph{initial marking} is
$\bfm_0 = \sum_{P \in \mathcal{P}} \mset{(P, c) \mid c \in \mathcal{I}(P)}$, where
every colour~$c$ is taken with its multiplicity in~$\mathcal{I}(P)$.

\paragraph{Bindings.}
A \emph{binding} of a transition~$T$ in a marking~$\bfm$ is a tuple
$\cvec = (c_1, \dots, c_m) \in \Dom^{km}$ such that
$\mathit{in}_T(\cvec) = \mset{(P(A_1), c_1), \dots, (P(A_m), c_m)} \le \bfm$, that is,
$\bfm$ contains a distinct token for every input arc. A binding~$\cvec$
\emph{enables}~$T$ if $T \in \mathcal{T}_I$, or if $T \in \mathcal{T}_G$ and
$\mathcal{G}_T(\cvec) = \mathit{True}$. The tokens produced by firing~$T$ on~$\cvec$
are $\mathit{out}_T(\cvec) = \mset{(P(A'_1), \pi_T^1(\cvec)), \dots,
(P(A'_n), \pi_T^n(\cvec))}$.

\paragraph{Firings.}
Within a round, we split the marking into two multisets: the multiset~$\bfm$
of the tokens of the snapshot that have not been consumed yet, and the
multiset~$\bfDelta$ of the tokens produced so far. Transitions consume tokens
only from~$\bfm$.

\begin{definition}[Firing]\label{def:firing}
  For multisets $\bfm$, $\bfDelta$, $\bfm'$, $\bfDelta'$, a
  transition~$T$, and a tuple $\cvec \in \Dom^{km}$, we write
  $\langle \bfm, \bfDelta \rangle \xrightarrow{T, \cvec}
    \langle \bfm', \bfDelta' \rangle$ iff:
  \begin{enumerate}
    \item $\cvec$ is a binding of~$T$ in~$\bfm$ that enables~$T$;
    \item if $T \in \mathcal{T}_G$, then no immediate transition has a
      binding in~$\bfm$ (rule~R0); and
    \item $\bfm' = \bfm - \mathit{in}_T(\cvec)$ and
      $\bfDelta' = \bfDelta + \mathit{out}_T(\cvec)$.
  \end{enumerate}
  We write $\rightarrow$ for the union of the relations
  $\xrightarrow{T, \cvec}$ over all $T$ and~$\cvec$, and $\rightarrow^*$ for
  its reflexive and transitive closure.
\end{definition}

Firing takes no time. Hence, the tokens in~$\bfm$ have the same colours as in
the snapshot, and guards are evaluated against the snapshot. The tokens
produced in a round are in~$\bfDelta$, so they cannot be consumed in the same
round, and every token of the snapshot is consumed at most once.

\begin{definition}[Round]\label{def:round}
  For markings $\bfm$ and~$\bfm'$, we write $\bfm \Rightarrow \bfm'$ iff
  there are multisets $\bfm''$ and~$\bfDelta''$ such that
  $\langle \bfm, \mset{} \rangle \rightarrow^* \langle \bfm'', \bfDelta''
  \rangle$, no firing is possible from
  $\langle \bfm'', \bfDelta'' \rangle$, and
  $\bfm' = \bfm'' + \bfDelta''$.
\end{definition}

That is, a round fires transitions until none is enabled by the remaining
tokens of the snapshot. Since every transition has an input arc, every firing
removes at least one token from the finite first component, so every sequence
of firings is finite. Hence, for every
marking~$\bfm$, there is a marking~$\bfm'$ with $\bfm \Rightarrow \bfm'$.

\begin{definition}[Clock tick]\label{def:tick}
  For markings $\bfm$ and~$\bfm'$ and a duration $h \in \Dom$ with $h > 0$,
  we write $\bfm \tickto{h} \bfm'$ iff $c + h \cdot \mathcal{V}_P(c) \in
  \Dom^k$ for every token $(P, c) \in \bfm$, and
  \[
    \bfm' = \mset{(P, c + h \cdot \mathcal{V}_P(c)) \mid (P, c) \in \bfm},
  \]
  where every token is taken with its multiplicity in~$\bfm$.
\end{definition}

That is, a tick of duration~$h$ moves the colour of every token by one Euler
step of size~$h$ along its rate of change. If $\Dom = \NatZero$, a tick is only
possible if no colour leaves~$\NatZero$.

\begin{remark}[Euler steps and exact flows]\label{rem:euler}
  In~\cite{EverdijBlom2005}, the colour of a token in place~$P$ evolves as the
  solution of the ordinary differential equation $\dot{c} = \mathcal{V}_P(c)$.
  The proper way to advance the colour by a duration~$h$ is therefore the exact
  flow $c \mapsto \Phi_P(c, h)$, where $\Phi_P(c, h)$ is the value at time~$h$
  of the solution that starts in~$c$. This requires $\Dom = \Real$ and the
  solutions to exist for all times. The Euler step $c + h \cdot
  \mathcal{V}_P(c)$ coincides with the exact flow if $\mathcal{V}_P$ is
  constant. Otherwise, it is an approximation with an error in $O(h^2)$ per
  tick. We use the Euler step, since the exact flow is, in general, neither
  available in closed form nor expressible in linear arithmetic.
\end{remark}

\begin{definition}[Execution]\label{def:execution}
  An \emph{execution} is an infinite sequence of markings
  $\bfm_0, \bfm_1, \bfm_2, \dots$ that starts with the initial
  marking~$\bfm_0$, together with an infinite sequence of clock values
  $0 = t_0 < t_1 < t_2 < \dots$ in~$\Dom$, such that for every
  $i \in \NatZero$, we have
  $\bfm_i \Rightarrow \bfm' \tickto{t_{i+1} - t_i} \bfm_{i + 1}$ for some
  marking~$\bfm'$. The marking~$\bfm_i$ is the marking at the clock
  value~$t_i$.
\end{definition}

The rules R0--R4 of~\cite{EverdijBlom2005} are reflected as follows. Rule~R0
is condition~2 of Definition~\ref{def:firing}. Rules R1 and~R3 follow from
maximality: all bindings in the snapshot that enable transitions and share no tokens
fire in the same round, that is, at the same clock value. Rules R2 and~R4
follow from consuming each token of the snapshot at most once: of several bindings
that share a token, only one fires. Everdij and Blom select it randomly,
whereas we select it nondeterministically.

\begin{remark}
  Our semantics deviates from~\cite{EverdijBlom2005} in two respects. First,
  tokens produced at the clock value~$t_i$ can only be consumed in the round
  at~$t_{i+1}$. In contrast, in~\cite{EverdijBlom2005}, an immediate
  transition enabled by freshly produced tokens fires at the same time. Hence,
  every stage of a cascade of immediate transitions costs one clock tick in
  our semantics. In return, rounds are finite. Second, the
  tokens produced in a round evolve during the following tick, that is, every
  token in~$\bfm_{i + 1}$ has evolved for the duration $t_{i+1} - t_i$,
  regardless of whether it was produced at~$t_i$.
\end{remark}

%%%%%%%%%%%%%%%%%%%%%%%%%
\bibliographystyle{plain}
\bibliography{references}

\end{document}
